\begin{abstract}
Direct-to-satellite (D2S) IoT devices often operate under severely constrained
knowledge of timing, frequency, and location.  Low-cost oscillators and sparse
satellite visibility make reliable uplink state estimation challenging, yet
coarse knowledge of device position and timing offset can substantially improve
uplink scheduling and link-budget management in non-terrestrial network (NTN)
scenarios.

This paper presents LEO-DTF: a Doppler-Time Fingerprint approach that exploits
the structured Doppler dynamics induced by low-earth-orbit (LEO) satellite
motion.  Using a two-line element (TLE) set and the SGP4 propagator, LEO-DTF
computes the expected range-rate and Doppler curve for any candidate position
inside a bounded region of interest (ROI).  A Bayesian grid-based posterior
estimator jointly accounts for unknown carrier frequency offset (CFO), oscillator
drift, and timing offset as nuisance parameters, yielding a maximum \emph{a posteriori}
(MAP) coarse location estimate together with uplink state quantities.
An ablation study in simulation demonstrates that multi-epoch Doppler-time
fingerprints reduce posterior ambiguity relative to single-snapshot Doppler,
even when nuisance parameters are present.

The current evaluation pipeline is entirely synthetic and trace-driven: it
provides a deterministic synthetic-pass generator, Monte Carlo coarse-localization
diagnostics, CRLB timestamp sensitivity analysis, a Doppler-time ambiguity
ablation study, and automatic table and figure generation.  A trace-driven
hardware-in-the-loop (HIL) validation plan is described using an LR1121/STM32
packet source (transmit-only) and a USRP B210 IQ capture front-end, with
offline Doppler/CFO/delay injection.  \textbf{No real satellite over-the-air
validation is claimed.}  Meter-level localization accuracy and GNSS replacement
are explicitly outside the scope of the current work.  Rather than replacing GNSS,
LEO-DTF is intended as a complementary coarse localization and uplink state-estimation
primitive for bounded-region D2S IoT scenarios, with current results read as
preliminary synthetic diagnostics.
\end{abstract}